\scanpage{302}
\mspage{452}

Notons que (33), s'écrit $\rho=-(1+N_1)\sigma$ i.e.
\deleted{\illegible{1}}
\[
\rho+\sigma=-N_1\sigma \tag{35}
\]
multipliant par $N_1$ à gauche on obtient
\[
N_1\rho+N_1\sigma=0
\]
(puisque $N_1^2=0$), d'ailleurs
\[
N_1\rho=\rho N_0,\qquad N_1\sigma=\sigma N_0
\]
(car $N_1=\rho N_0\rho^{-1}=\sigma N_0\sigma^{-1}$), donc (\unsure{35}) s'écrit
\[
\rho+\sigma=-N_1\sigma=+N_1\rho=-\sigma N_0=+\rho N_0
\overset{\text{déf}}{=}N . \tag{\unsure{35|36}}
\]
d'ailleurs quand $\rho$ est sous la forme générale
$\rho=\begin{pmatrix}0&u\\v&w\end{pmatrix}$ on trouve
\[
N\deleted{\illegible{1}}=\begin{pmatrix}0&0\\0&u\end{pmatrix},
\qquad N^2=\begin{pmatrix}0&0\\0&u^2\end{pmatrix}
\]
donc la normalisation (34) $u=1$ signifie aussi
\[
\rho+\sigma=\begin{pmatrix}0&0\\0&1\end{pmatrix}\deleted{\unsure{par}}, \tag{37}
\]
\inserted{ou encore (comme $u$ est inversible)
\[
N^2=N\qquad (\text{où }N=\rho+\sigma) \tag{37 bis}
\]}
ou encore que $\sigma$ est donné en termes de $\rho$ par
\[
\sigma=\begin{pmatrix}0&-1\\-v&1-w\end{pmatrix}. \tag{38}
\]
d'ailleurs, $N_0$ \inserted{$=\rho N_1\rho^{-1}$} sera de la forme
$\begin{pmatrix}0&\alpha\\0&0\end{pmatrix}$,
un calcul immédiat montre que $\alpha$ est donné par la formule
$\unsure{u}=\alpha v$, i.e. (\unsure{si} on normalise par $u=1$)

\margin{\uline{NB} On voit \unsure{aussitôt} que les conditions envisagées jusqu'à présent
\unsure{impliquent} que $1,\rho,\sigma$ \unsure{sont lin. indép.}
\unsure{ds. chaque} fibre, i.e. $T_\rho\cap T_\sigma=\unsure{\mathbb G_m}$ !
i.e. $\widetilde T_\rho\cap\widetilde T_\sigma=\{1\}$.}

\scanpage{303}

\[
N_0=\begin{pmatrix}0&1/v\\0&0\end{pmatrix},\qquad
\varepsilon_0=\begin{pmatrix}1&1/v\\0&1\end{pmatrix}. \tag{39}
\]

\margin{\uline{NB} On aura
\[
\boxed{\sigma^{-1}\rho=-\varepsilon_0}\deleted{\unsure{N_0}} \tag{39 bis}
\]
Car $\sigma^{-1}\rho=\sigma^{-1}(\rho\sigma^{-1})\sigma
=\sigma^{-1}(-\varepsilon_1)\sigma
=-\sigma^{-1}(\varepsilon_1)=-\varepsilon_0$.}

\deleted{Revenons}
Ces normalisations canoniques étant faites, revenons à la proposition de la page 436 ; comme $\varepsilon_1$ a été choisi
(de façon tt ce qu'il y a de canonique) il n'y a plus d'inconvénient de
\deleted{\illegible{1}}\inserted{laisser} l'équation (32) (mise sous la forme en apparence plus intrinsèque
(11) \unsure{ou} (12)) sous la forme \deleted{\illegible{1}} initiale,
où a priori on songeait à introduire le paramètre $\gamma$ avec
\[
[\alpha,\rho]=[\beta,\rho]\varepsilon_1(\gamma)\unsure{\,;}\quad -
\]
mais il apparaît que c'est $\gamma+1=\mu$ qui a une signification particulièrement remarquable
dans la théorie --- aussi en \unsure{gardons|garderons} la forme (32), avec le facteur $\varepsilon_1(\mu-1)$.

L'expression (97) de $P(\unsure{\zeta})$, grâce à (37), s'écrit
\[
P(\zeta)=D\qquad\text{pour }\zeta=\begin{pmatrix}A&B\\C&D\end{pmatrix}, \tag{40}
\]

\scanpage{304}
\mspage{453}

Donc la proposition nous dit que le schéma \inserted{$\mathcal X_{\rho,\sigma}$}
des solutions simultanées en $(\zeta,\eta,\mu)$ de
\[
\det\zeta=\det\eta=1,\qquad \zeta=\eta\varepsilon_1(\mu-1),\qquad
\operatorname{Tr}(\zeta\rho)=\operatorname{Tr}\rho,\qquad
\operatorname{Tr}\eta\deleted{\unsure{\zeta|\rho}}\inserted{\sigma}
=\operatorname{Tr}\unsure{\sigma|\alpha} \tag{41}
\]
est isomorphe, via $(\zeta,\eta,\mu)\longrightarrow\zeta$, à celui des
$\zeta=\begin{pmatrix}A&B\\C&D\end{pmatrix}$ satisfaisant les conditions
\[
\underbrace{\det\zeta}_{AD-BC}=1,\qquad
\underbrace{\operatorname{Tr}\zeta\rho}_{(C+wD)+vB}
=\underbrace{\operatorname{Tr}\rho}_{w},\qquad
\text{et }D\text{ section inversible de }\mathcal O_S \tag{42}
\]
et on retrouve $\mu$ en termes de $\zeta$ par
\[
\mu=1/D \tag{43}
\]
(ici $c=\operatorname{Tr}(\rho+\sigma)=\operatorname{Tr}\begin{pmatrix}0&0\\0&1\end{pmatrix}=1$),
d'où $\eta$ par $\eta=\zeta\varepsilon_1(1-\mu)$.

Le sous-schéma en groupes $B'_0$ de la page 442$'$ n'est autre que le schéma des matrices
\[
B'_0=\left\{\begin{pmatrix}\mu&\lambda\\0&1\end{pmatrix}
\,\middle|\,\mu\in\Gamma\mathbb G_m,\ \lambda\in\Gamma\mathbb G_a\right\} \tag{44}
\]
Nous allons expliciter l'isomorphisme (79)
\[
\underline u\longmapsto\bigl([\underline u,\rho],[\underline u,\sigma],\det\underline u\bigr) \tag{45}
\]
en termes des coefficients $\mu,\lambda$ de $B'_0$, à l'aide des paramètres
$v,w$ dans $\rho=\begin{pmatrix}0&1\\v&w\end{pmatrix}$.

\scanpage{305}

On trouve
\deleted{\unsure{$\underline u(\rho)=$}}
\[
\begin{aligned}
{}[\underline u,\rho]
&=\begin{pmatrix}\mu&\lambda\\0&1\end{pmatrix}
\begin{pmatrix}0&1\\v&w\end{pmatrix}
\frac1\mu\begin{pmatrix}1&-\lambda\\0&\mu\end{pmatrix}
\left(-\frac1v\right)\begin{pmatrix}w&-1\\-v&0\end{pmatrix}\\
&=-\frac1{\mu v}\begin{pmatrix}\lambda v&\mu+\lambda w\\v&w\end{pmatrix}
\begin{pmatrix}w+\lambda v&-1\\-\mu v&0\end{pmatrix}\\
&=-\frac1{\mu v}
\begin{pmatrix}
\deleted{\illegible{2}\unsure{-\mu)vw+\lambda}}\inserted{v^2\lambda^2-v\mu^2-vw\lambda\mu+vw\lambda}&-\lambda v\\
v^2\lambda-vw\mu+vw&-v
\end{pmatrix}\\
&=\frac1\mu\begin{pmatrix}
-v\lambda^2+w\lambda\mu+\mu^2-w\lambda&\unsure{\lambda}\\
-v\lambda+w\mu-w&1
\end{pmatrix}
\end{aligned}
\]

\margin{\unsure{or|et} pour $\sigma=\begin{pmatrix}0&-1\\v'&w'\end{pmatrix}$
avec $v'=-v$, $w'=1-w$}

\[
\begin{aligned}
{}[\underline u,\sigma]
&=\begin{pmatrix}\mu&\lambda\\0&1\end{pmatrix}
\begin{pmatrix}0&-1\\v'&w'\end{pmatrix}
\frac1\mu\begin{pmatrix}1&-\lambda\\0&\mu\end{pmatrix}
\left(\frac1{v'}\right)\begin{pmatrix}w'&1\\-v'&0\end{pmatrix}\\
&=\frac1{\mu v'}\begin{pmatrix}\lambda v'&-\mu+\lambda w'\\v'&w'\end{pmatrix}
\begin{pmatrix}w'+\lambda v'&1\\-\mu w'&0\end{pmatrix}\\
&=\frac1{\mu v'}\begin{pmatrix}
v'^2\lambda^2+v'\mu^2-v'w'\lambda\mu+v'w'\lambda&\lambda v'\\
v'^2\lambda-v'w'\mu+v'w'&v'
\end{pmatrix}\\
&=\frac1\mu\begin{pmatrix}
v'\lambda^2\unsure{-}\deleted{\unsure{\mu}}w'\lambda\mu+\mu^2+w'\lambda&\lambda\\
v'\lambda-w'\mu+w'&1
\end{pmatrix}\\
&=\frac1\mu\begin{pmatrix}
-v\lambda^2+(w-1)\lambda\mu+\mu^2\unsure{-}(w-1)\lambda&\lambda\\
-v\lambda+(w-1)\mu-(w-1)\deleted{\unsure{\lambda}}&1
\end{pmatrix}
\end{aligned}
\]

En résumé
\[
\left\{\begin{aligned}
\zeta=[\underline u,\rho]
&=\frac1\mu\begin{pmatrix}
\mu^2+(-v\lambda+w(\mu\unsure{-}1))\lambda&\lambda\\
-\deleted{\unsure{(}}v\lambda+w(\mu-1)&1
\end{pmatrix}\\
\eta=[\underline u,\sigma]
&=\frac1\mu\begin{pmatrix}
\mu^2+(-v\lambda+(w-1)(\mu-1))\lambda&\lambda\\
-v\lambda+(w-1)(\mu-1)&1
\end{pmatrix}\quad\text{et, bien sûr}\\
\mu&=\det\underline u\ \deleted{\unsure{=\mu}}
\end{aligned}\right. \tag{46}
\]
donc on a, si $\zeta=\begin{pmatrix}A&B\\C&D\end{pmatrix}$
\[
B=\lambda/\mu,\quad D=1/\mu\qquad\text{i.e.} \tag{47}
\]
\[
\mu=1/D,\quad\lambda=B/D,\qquad\text{\unsure{soit|dont} }\underline u=\begin{pmatrix}1/D&B/D\\0&1\end{pmatrix} \tag{48}
\]

\scanpage{306}
\mspage{453bis}

De même, si $\eta=\begin{pmatrix}R&S\\T&U\end{pmatrix}$, on aura \deleted{$S=B,U=D$}
\deleted{d'où}
\[
\left\{\begin{aligned}
S=B&=\lambda/\mu,\quad U=D=1/\mu\qquad\text{d'où}\\
\deleted{\unsure{et}}\ \mu&=1/U,\quad\lambda=S/U
\end{aligned}\right. \tag{47 bis}
\]
donc
\[
\underline u=\begin{pmatrix}1/U\deleted{\unsure{(=\mu)}}&S/U\\0&1\end{pmatrix}. \tag{\unsure{48 bis}}
\]

Considérons maintenant l'équation
\[
[\alpha,\rho]=[\beta,\sigma]\varepsilon_1(\mu-1),\qquad\text{avec} \tag{49}
\]
\[
\alpha=\begin{pmatrix}a&b\\c&d\end{pmatrix},\quad
\beta=\begin{pmatrix}r&s\\t&u\end{pmatrix}, \tag{\unsure{50}}
\]
et posons, pour des calculs plus \unsure{élégants}
\[
\rho=\begin{pmatrix}0&\unsure{v_1}\\\unsure{v}&w\end{pmatrix},\quad
\sigma=\begin{pmatrix}0&v'_1\\v'&w'\end{pmatrix} \tag{51}
\]
(on a écrit $v_1,v'_1$ au lieu de $u,u'$, pour ne pas \unsure{interférer} avec la notation $u$ pour un des coefficients de $\beta$), avec
\[
v'=-v,\quad v'_1=-\unsure{v_1|v'},\quad w+w'=v_1,\qquad\text{posons} \tag{52}
\]
\[
\delta_\alpha=\det\alpha=ad-bc,\quad
\delta_\beta=\det\beta=ru-st,
\tag{53}
\]

\margin{On ne pose pas $v_1=1$, pour \unsure{n'avoir} pas à \unsure{refaire} le $\overline{\mathrm m}$ calcul pour $[\beta,\sigma]$ que pour $[\alpha,\rho]$ !}

on trouve alors
\[
\begin{aligned}
\zeta=[\alpha,\rho]
&=\begin{pmatrix}a&b\\c&d\end{pmatrix}
\begin{pmatrix}0&v_1\\v&w\end{pmatrix}
\frac1{\delta_\alpha}\begin{pmatrix}d&-b\\-c&a\end{pmatrix}
\deleted{\unsure{(}}\frac1{-vv_1}\begin{pmatrix}w&-v_1\\-v&0\end{pmatrix}\\
&=\unsure{-}\frac1{\delta_\alpha vv_1}
\begin{pmatrix}bv&av_1+bw\\dv&cv_1+dw\end{pmatrix}
\begin{pmatrix}dw+bv&-dv_1\\-cw-av&cv_1\end{pmatrix}
=\begin{pmatrix}A&B\\C&D\end{pmatrix}
\end{aligned}
\]
avec
\[
\left\{\begin{aligned}
(-\delta_\alpha vv_1)A&=\ldots\quad\text{(barbe).},
&(-\delta_\alpha vv_1)B&=-vv_1bd+v_1^2ac+v_1wbc\\
(-\delta_\alpha vv_1)C&=\ldots\quad\text{(barbe)},
&(-\delta_\alpha vv_1)D&=-vv_1d^2+v_1wcd+v_1^2c^2
\end{aligned}\right. \tag{54}
\]
\margin{NB les formules pour $B,D$ se \unsure{simplifient} par $v_1$}

\scanpage{307}

et on trouve de $\overline{\mathrm m}$
\[
\eta=\unsure{(\beta,\sigma)|[\beta,\sigma]}=
\begin{pmatrix}R&S\\T&U\end{pmatrix}
\]
avec
\[
\tag{55}
\begin{array}{ll}
(-vv_1\delta_\beta)R=\ldots
&(-vv_1\delta_\beta)S
 =(v_1v)su+v_1^{2}tr-v_1w'st\\
(-vv_1\delta_\beta)T=\ldots
&(-vv_1\delta_\beta)U
 =-(v_1v)u^2-v_1w'tu+v_1^{2}t^2.
\end{array}
\]
\margin{\uline{NB} on peut diviser par $v_1$.}

\unsure{On} en tire en particulier
\deleted{\illegible{1}}
\[
\tag{56}
\left\{
\begin{aligned}
D&=\frac{-v\deleted{\unsure{v_1}}d^2+
\deleted{\unsure{v_1}}wcd+v_1^{\deleted{2}}c^2}
{-v\deleted{\unsure{v_1}}(ad-bc)}\\
\mu=1/D&=\frac{-v\deleted{\unsure{v_1}}(ad-bc)}
{-v\deleted{\unsure{v_1}}d^2+
\deleted{\unsure{v_1}}wcd+v_1^{\deleted{2}}c^2}\\
\deleted{(57)}\quad\lambda&=\deleted{\unsure{C/D}}\inserted{B/D}
=\frac{-v\deleted{\unsure{v_1}}bd+v_1^{\deleted{2}}ac+
\deleted{\unsure{v_1}}wbc}
{-v\deleted{\unsure{v_1}}d^2+
\deleted{\unsure{v_1}}wcd+v_1^{\deleted{2}}c^2}.
\end{aligned}
\right.
\]
ce qui permet donc d'\unsure{expliciter}
\[
\underline{u}=\begin{pmatrix}1/D=\mu&B/D=\lambda\\0&1\end{pmatrix}
\]
en termes de $a,b,c,d$ (et des paramètres $u_1,v,w$
qui déterminent $\rho$).

Pour \uline{$v_1=1$}, la formule pour $D$ de (56), s'écrit
\[
\tag{57}
D=\det(d\rho+c)/-v\delta_\alpha
=\det(d\rho+c)/\det\rho\det\alpha
\]
\deleted{la condition \unsure{à} \illegible{1} sur $\alpha$ ``$D$ inv.''}
\[
\tag{58}
D(\alpha)\text{ inversible}\quad\Longleftrightarrow\quad
d\rho+c\text{ inversible}.
\]
mais je m'aperçois que c'est une forme \unsure{plate},

\scanpage{308}
\mspage{454}

\uline{V}\qquad
\uline{Étude}\inserted{\uline{(de l'équation)}}\uline{ dans}
$GP(1)$, \uline{et la condition} $\mathrm{Tr}\,\sigma=0$.

Jusqu'à présent nous avons étudié l'équation (32)
\[
\tag{*}
[\alpha,\rho]=[\beta,\unsure{\sigma}]\,
\unsure{\varepsilon_1(\mu-1)}
\]
dans $G\deleted{\unsure{\ell_2}}$ ($\simeq\mathrm{Gl}(2)$), ou ce qui revient au
$\overline{\mathrm m}$, \deleted{dans} en tant que relation dans $\mathrm{Sl}(2)$.
Ici $\alpha,\beta$ sont des sections \deleted{\illegible{1}}\inserted{variables} de
$G\deleted{\unsure{\ell_2}}\simeq\mathrm{Gl}(2)$, mais $[\alpha,\rho]$ et
$[\beta,\sigma]$ ne dépendent que des sections \unsure{correspondantes}
$\widetilde\alpha,\widetilde\beta$ de
\[
PG\ (\simeq GP(1))\simeq
G/\underbrace{\unsure{G_m}}_{\text{Centre de }G}
\simeq\mathrm{Aut}_{\unsure{1|l}}(G)
\]
(de $\overline{\mathrm m}$ qu'elles ne dépendent de $\rho,\sigma$ que
\unsure{via} $\widetilde\rho,\widetilde\sigma$).
On peut donc \unsure{considérer} que c'est une équation
\inserted{dans \deleted{$\mathrm{Sl}(2)$} $SG\simeq\mathrm{Sl}(2)$,}
liant $\widetilde\alpha,\widetilde\beta,\mu$.

On peut se proposer d'étudier l'équation correspondante dans
$\unsure{PG}\simeq GP(1)$,
\[
\tag{1}
[\widetilde\alpha,\widetilde\rho]
 =[\widetilde\beta,\widetilde\sigma]\cdot\widetilde\varepsilon_1(\mu-1)
\]
ce qui s'écrit aussi dans $\mathrm{Gl}(2)$, sous la forme \unsure{équivalente}
\[
\tag{2}
[\alpha,\rho]=\varepsilon[\beta,\sigma]\,
\overset{\deleted{\unsure{N_0}}}{\varepsilon_1}(\mu-1),
\]
où $\varepsilon$ est une section de $\unsure{G_m}$,
\deleted{et en}\inserted{bien} \deleted{fait} entendu uniquement
\unsure{déterminé|déterminée} par $\alpha,\beta,\mu$,
(et $\overline{\mathrm m}$ par $\widetilde\alpha,\widetilde\beta,\mu$)).
On aura nécessairement

\scanpage{309}

\[
\tag{3}
\varepsilon^2=1
\]
(prendre le dét. des deux \unsure{mbs}). On considère (2) comme une équation
reliant les variables $\alpha,\beta,\mu,\varepsilon$
($\alpha,\beta$ dans $\mathrm{Gl}(2)$, $\mu,\varepsilon$ dans
$\unsure{\mathbb E^1|\mathbb E'}$ -- et en fait nécessairement dans
$\unsure{\mathbb G_m|G_m}$, et $\varepsilon$ \unsure{déjà} dans
$\mu_2={}_2\unsure{\mathbb G_m|G_m}$), ou aussi
$\widetilde\alpha,\widetilde\beta,\mu,\varepsilon$.
Posons, \deleted{comme dans}\inserted{\unsure{en analogie avec}} (II),
\[
\tag{4}
\zeta=[\alpha,\rho],\qquad
\eta=\deleted{\illegible{1}}\varepsilon
[\unsure{\uuline{\beta}},\sigma]\deleted{\unsure{\varepsilon_1 t|\text{et}}},
\]
de sorte que (2) prend la forme
\[
\tag{5}
\zeta=\eta\varepsilon_1(\mu-1)
\]
comme dans (II). Pour pouvoir appliquer (II) à cette situation, il faudra
supposer que $\eta=\varepsilon[\beta,\sigma]$ satisfait aux
$\overline{\mathrm m}^{s}$ conditions préliminaires que
\[
\eta_0=[\beta,\sigma]
\]
lui $\overline{\mathrm m}$, savoir $\det\eta=1$
($\Longleftrightarrow\ \varepsilon^2=1$) et
\deleted{\unsure{que}}
\[
\underbrace{\mathrm{Tr}\,\eta}_{\inserted{\varepsilon\mathrm{Tr}\,\eta_0
=\varepsilon\mathrm{Tr}\,\sigma}}
=\mathrm{Tr}\,\sigma
\qquad\deleted{\unsure{car}\ \varepsilon\mathrm{Tr}\,\eta_0=}
\]
i.e.
\[
\varepsilon\mathrm{Tr}\,\sigma=\mathrm{Tr}\,\sigma.
\]

\scanpage{310}
\mspage{455}

\uline{Lemme.} Soit $\sigma$ une section de $G$ (forme de
$\mathrm{Gl}(2)$, \unsure{associé|associée} à une Alg d'Azumaya
$\unsure{\mathcal M|M}$) -- i.e. une section de
$\unsure{\mathcal M|M}$ telle que $\det\sigma$ inversible.
\deleted{Conditions équivalentes}\inserted{Soit $A_\sigma$ le schéma}
des sections $\eta$ \deleted{\illegible{1}} de $G$ satisfaisant les conditions
\[
\det\eta=1\quad\text{(i.e. $\eta$ dans $SG$)},\qquad
\mathrm{Tr}\,\eta\circ\sigma=\mathrm{Tr}\,\sigma.
\]
Conditions équivalentes a) b) c) :

a) $A_\sigma$ est stable par translations par
$\mu_2\subset\unsure{\mathbb G_m|G_m}$ :
\[
\unsure{\eta}\longmapsto\varepsilon\eta\qquad(\varepsilon^2=1)
\]

b) $\mathrm{Tr}\,\sigma=0$

\deleted{c) $\sigma^2\in\Gamma\unsure{\mathbb G_m|G_m}$ i.e.
$\widetilde\sigma^2=1$.}

\unsure{d)|c)} (Dans les \unsure{conditions} du n° précédent
\inserted{\uline{$\overline{\mathrm{IV}}$} (p. 452)},
avec les normalisations (34) (38) pour $\rho,\sigma$) On a
\[
\tag{6}
\sigma=\begin{pmatrix}0&-1\\v'&0\end{pmatrix}
=\begin{pmatrix}0&-1\\-v&0\end{pmatrix}
\]
i.e. (avec les notations du loc. cit.) on a $w=1$,
\[
\tag{7}
\text{i.e.}\qquad
\rho=\begin{pmatrix}0&1\\v&1\end{pmatrix}.
\]

\inserted{d) On a $\sigma^2\in\Gamma\unsure{\mathbb G_m|G_m}$,
i.e. $\widetilde\sigma^2=1$}

\inserted{Enfin, ces conditions impliquent la suivante d), et
\unsure{lui|leur} sont équivalentes \unsure{si $\sigma$}
\deleted{\unsure{quand $\sigma$}} n'est une homothétie sur
aucune \unsure{fibre} (\unsure{donc} i.e. $\sigma$ est régulier au sens de
Steinberg).}

\deleted{Ceci signifie que}\inserted{\unsure{L'équivalence}
$a)\Longleftrightarrow b)$ signifie \unsure{aussi|ainsi}}
que $\unsure{\mathbb E^1|\mathbb E'}{}^{\mu_2}=(0)$, qui est bien
\unsure{connu} et de vérification immédiate.
L'équivalence de b) et d) est immédiate, reste à prouver
$b)\Longleftrightarrow c)$.
\deleted{\unsure{Supposons|Supposant} cette condition \unsure{satisfaite}}
\inserted{Or on a l'équation \unsure{de|d'}
\unsure{Hamilt-Cayley|Hamilton-Cayley}}
\[
\sigma^2-(\mathrm{Tr}\,\sigma)\sigma+\det\sigma=0
\]

\scanpage{311}

qui montre que $\mathrm{Tr}\,\sigma=0\Longrightarrow\sigma=\det\sigma\cdot1$,
et inversement si $\sigma^2$ est un scalaire, on aura que
$(\mathrm{Tr}\,\sigma)\sigma=\sigma^2-\det\sigma\cdot1$ en est un,
ce qui n'est possible (si $\sigma$ n'est \deleted{\unsure{pas}} un scalaire sur aucune
fibre) que si $\mathrm{Tr}\,\sigma=0$.

Supposons pour la suite ces conditions satisfaites,
\deleted{on \unsure{trouve} alors que le schéma $\mathcal G_{\rho,\sigma}$
des \illegible{1} $(\alpha,\beta,\mu,\varepsilon)$ satisfaisant}
\inserted{Soit}
\[
\tag{8}
\mathcal G^*_{\rho,\sigma}\subset
G\times G\times\unsure{\mathbb E^1|\mathbb E'}
\times\unsure{\mathbb E^1|\mathbb E'}
\]
le schéma des solutions en $(\alpha,\beta,\mu,\varepsilon)$ de l'équation (2),
qui est aussi l'image inverse dans
$G\times G\times\unsure{\mathbb E^1|\mathbb E'}
\times\unsure{\mathbb E^1|\mathbb E'}$ \unsure{d'un} sous schéma
\deleted{$\widetilde{\mathcal G}_{\rho,\unsure{\sigma}}$
\illegible{1} $PG\times$}
\[
\tag{9}
\widetilde{\mathcal G}^*_{\rho,\sigma}\subset
PG\times PG\times\unsure{\mathbb E^1|\mathbb E'}
\deleted{\times\unsure{\mathbb E^1|\mathbb E'}}
\]
formé des solutions \inserted{$(\alpha,\beta,\mu)$} de (1).
La projection sur le \unsure{deuxième|dernier} facteur
$\mathbb E^1$ fournit un morphisme
\[
\tag{10}
\mathcal G^*_{\rho,\sigma}\xrightarrow{\underline{\varepsilon}}\mu_2,
\]

\scanpage{312}
\mspage{456}
qui se \unsure{factorise} d'ailleurs par $\widetilde{\mathcal G}^{*}_{\rho,\sigma}$ :
\[
\tag{11}
\widetilde{\varepsilon}:\widetilde{\mathcal G}^{*}_{\rho,\sigma}\longrightarrow\mu_2.
\]
Notons qu'on a
\[
\deleted{\unsure{(11)}}\qquad
\mathcal G_{\rho,\sigma}\subset\mathcal G^{*}_{\rho,\sigma}
\]
plus précisément
\[
\tag{12}
\mathcal G_{\rho,\sigma}=\underline{\varepsilon}^{-1}(1)
\]
ce qui conduit à introduire aussi
\[
\tag{13}
\underset{\substack{\Vert\ \text{déf}\\
\underline{\widetilde{\varepsilon}}^{-1}(1)}}{\widetilde{\mathcal G}_{\rho,\sigma}}
=\widetilde{\mathcal G}^{*}_{\rho,\sigma},
\]
de sorte que $\mathcal G_{\rho,\sigma}$ est l'image inverse dans
$G\times G\times\unsure{\mathbb E^1|\mathbb E'}$ de $\widetilde{\mathcal G}_{\rho,\sigma}$.

On a un morphisme
\[
\tag{14}
\left\{\begin{aligned}
(\alpha,\beta,\mu,\varepsilon)&\longrightarrow
\left(\underbrace{(\zeta=\alpha\rho(\alpha)^{-1},\,
\eta=\varepsilon\beta\sigma(\beta)^{-1},\,\mu)}_{\in\Gamma\mathcal X_{\rho,\sigma}},\,
\underbrace{\varepsilon}_{\in\Gamma\mu_2}\right)\\
\text{de}\quad\mathcal G^{*}_{\rho,\sigma}&\longrightarrow\mathcal X_{\rho,\sigma}\times\mu_2
\end{aligned}\right.
\]
et l'image inverse d'une section $((\zeta,\eta,\mu),\varepsilon)$ du second membre est le
schéma des solutions en $\alpha,\beta$ des équations
\[
\tag{15}
{}^{(a)}\underset{\inserted{\unsure{\alpha}}}{\alpha\rho(\alpha)^{-1}}=\zeta,
\qquad {}^{(b)}\beta\sigma(\beta)^{-1}=\varepsilon\eta
\]

\scanpage{313}
\unsure{Or|On} on a, via l'isom
\[
B'_0\xrightarrow{\sim}\mathcal X_{\rho,\sigma}
\]
une solution privilégiée $\alpha=\underline u$ de la première équation (15), qui est
\unsure{caractérisée} comme l'unique solution qui soit \deleted{\unsure{une}} section de
$B'_0$ (\unsure{unicité} résultant de
\[
T_\alpha\cap B'_0=\{1\}
\]
).
Mais on \unsure{aura}, pour cet $u$
\[
\underbrace{u\sigma(u)^{-1}}_{[u,\sigma]}=\eta
\qquad(\text{et non }\varepsilon\eta\ !)
\]
pour trouver en termes de $u$ une solution de l'équation
\inserted{(15 b)} $\beta\sigma(\beta)^{-1}=\varepsilon\eta$, on est conduit à essayer sous la forme
\[
\tag{16}
\beta=u\tau\quad\text{d'où}\quad
\beta\deleted{\unsure{\rho}}\inserted{\sigma}(\beta)^{-1}
=\deleted{\unsure{\beta}}\inserted{u}\tau\sigma(\tau)^{-1}
\sigma(\deleted{\unsure{\beta}}\inserted{u})^{-1},
\]
en essayant de déterminer $\tau$ de façon qu'on ait
\[
\tag{17}
\underbrace{\tau\sigma(\tau)^{-1}}_{[\tau,\sigma]}=\varepsilon,
\]
ce qui implique bien
\[
\beta\sigma(\beta)^{-1}
=\deleted{\unsure{\beta}}\inserted{u}\varepsilon\sigma(\deleted{\unsure{\beta}}\inserted{u}^{-1})
=\varepsilon\deleted{\unsure{\beta}}\inserted{u}\sigma(\deleted{\unsure{\beta}}\inserted{u})^{-1}
=\varepsilon\eta
\]
i.e. (15 b). \deleted{\unsure{Lemme}}

\scanpage{314}
\mspage{457}
La relation (17) s'écrit aussi
\[
\deleted{\sigma(\tau)=\varepsilon\tau}
\qquad(\text{prendre l'inverse})
\]
\[
\tag{18}
\tau\sigma\tau^{-1}\sigma^{-1}=\varepsilon
\quad\text{ou}\quad
\sigma\tau\sigma^{-1}\tau^{-1}=\varepsilon
\quad\text{ou}
\]
\[
\tag{19}
\sigma=\varepsilon\tau(\sigma)
\quad\text{ou}\quad
\tau(\sigma)=\varepsilon\sigma
\]
ce qui implique que $\tau$ \uline{normalise} $T_\sigma$
\[
\tag{20}
\tau(T_\sigma)=T_\sigma\quad\text{i.e.}\quad\tau\in\Gamma N_\sigma
\]
où
\[
\tag{21}
N_\sigma\overset{\text{déf}}{=}\operatorname{Norm}_G(T_\sigma).
\]

\uline{Corollaire au lemme} (p.455) les conditions \inserted{\unsure{équiv.}} du lemme
\deleted{\unsure{équivalent}} \inserted{sont impliquées}
\deleted{aux} \inserted{par les} suivantes ;
\inserted{et \unsure{l'impliquent} si $\sigma$ est \unsure{régulière|régulier} :}

(e) $\forall$ section \inserted{locale fppf}
$\varepsilon$ de $\mu_2$, existe \inserted{\unsure{loc\textsuperscript{t}|loc.}}
section \deleted{\illegible{1}} $\tau$ de $G$ telle qu'on ait (17) (ou encore (19))

(f) (Si $2$ inv. sur $S$) idem, avec $\varepsilon=-1$.

Si on a (e), alors on voit que pour tout \deleted{\illegible{1}} $\varepsilon$,
on doit avoir $\operatorname{Tr}\varepsilon\sigma=\operatorname{Tr}\sigma$
or $\operatorname{Tr}\varepsilon\sigma=\varepsilon\operatorname{Tr}\sigma$, donc
\inserted{(faisant varier $\varepsilon$)}
$\operatorname{Tr}\sigma\in\Gamma(\unsure{\mathbb E^1|\mathbb E'}^{\mu_2}=0)$
donc $\operatorname{Tr}\sigma=0$. Si on le suppose seulement pour $\varepsilon=-1$,
on aura $\operatorname{Tr}\sigma=-\operatorname{Tr}\sigma$ i.e.
$2\operatorname{Tr}\sigma=0$, i.e. $\operatorname{Tr}\sigma=0$ si $2$ inv.
Inversement, supposons $\operatorname{Tr}\sigma=0$, alors
$\operatorname{Tr}\varepsilon\sigma=0$, d'ailleurs
$\det(\varepsilon\deleted{\unsure{\sigma}})=\varepsilon^2=1$, donc $\varepsilon$ satisfait
aux ``conditions numériques'' pour être de la forme $[\tau,\sigma]$ pour
$\tau$ convenable. Pour

\scanpage{315}
\unsure{assurer} l'existence locale de $\tau$, il reste à s'assurer que
$\varepsilon\sigma$ est section régulière au sens de Steinberg i.e. n'est nulle part
un scalaire, ce qui provient du fait qu'il en est ainsi de $\sigma$.

Ainsi, sous les conditions actuelles, soit $N'_\sigma\subset N_\sigma$ le sous-schéma de
\deleted{$N_\sigma$} $G$ (ou de $\unsure{N_\sigma|N}$) défini par
\[
\tag{22}
N'_\sigma=\{\tau\mid[\tau,\sigma]\in\Gamma\mu_2\}
\]
alors $N'_\sigma$ est stable par translation à droite par $T_\sigma$ --- mieux,
c'est un sous-schéma en groupes de $N_\sigma$ \inserted{(contenant $T_\sigma$)}, car si
$\tau,\tau'$ sont tels que
\[
\tau(\sigma)=\varepsilon\sigma,\qquad\tau'\sigma=\varepsilon'\sigma,
\]
on aura
\[
\tau'\tau(\sigma)=\varepsilon'\varepsilon\cdot\sigma,
\]
en $\overline{\mathrm m}$ temps on voit que
\[
\tau\longmapsto\varepsilon
\]
est un hom de groupes $N'_\sigma\longrightarrow\mu_2$, et on aura une suite exacte de
schémas en groupes
\[
\tag{23}
1\longrightarrow T_\sigma\hookrightarrow N'_\sigma
\xrightarrow{\varepsilon_\sigma}\mu_2\longrightarrow1.
\]
On \unsure{notera} que si $\tau\in\Gamma G$ normalise $T_\sigma$, on aura
\[
\tau(\sigma)=a\sigma+b\qquad a,b\in\Gamma\mathcal O_S,
\]

\scanpage{316}
\mspage{\unsure{458|488}}
et prenant les traces des deux mbs on trouve (compte tenu que
$\operatorname{Tr}\tau(\sigma)=\operatorname{Tr}\sigma=0$)
$\operatorname{Tr}b.1=0$ i.e. $2b=0$, donc si $b$ est inv.
$b=0$, donc $\tau(\sigma)=a\sigma$ et on a nécessairement $a^2=1$ i.e. $a$ est un
$\varepsilon$. Donc
\[
\tag{24}
N'_\sigma=N_\sigma\quad\text{en car. résiduelle }\ne2.
\]

\deleted{\unsure{Mais en car}}
Plus \unsure{généralement}, sur une base $S$ telle que
$2$ \deleted{ne} soit \deleted{\unsure{pas diviseur de zéro}} \inserted{section régulière}
(i.e. non div. de zéro) dans $\Gamma(S,\mathcal O_S)$, on voit que toute section de
$\unsure{N'_\sigma|N_\sigma}$ est en fait une section de $N_\sigma$.
Mais il n'en est plus ainsi en car $2$, car
\[
\sigma'=a\sigma+b\qquad(a\in\Gamma G_m,\ b\in\Gamma\mathcal O_S)
\]
dès qu'il satisfait
\[
\operatorname{Tr}\sigma'=\unsure{\sigma|0},\quad\det\sigma'=\det\sigma
\]
est conjugué \unsure{loc\textsuperscript{t}} à $\sigma$ --- et
\deleted{\unsure{il suffit pour}} \inserted{les conditions}
s'écrivent (comme
$\unsure{\det}(a\sigma+b)=a^2\det\sigma+
 ab\underbrace{\deleted{\unsure{\det}}\inserted{\operatorname{Tr}}\sigma}_{0}+b^2$
\[
\left\{\begin{aligned}
-va^2+b^2&=-v\\
2b&=0\qquad(\text{automatique en car. }2)
\end{aligned}\right.
\]
qui n'impliquent aucunement $b=0$ ....

\deleted{\illegible{1}} Considérons maintenant le groupe produit
\[
B'_0\times A_\rho\times N'_\sigma,
\]

\scanpage{317}

et le morphisme de schémas
\[
\tag{25}
\left\{\begin{aligned}
B'_0\times A_\rho\times N'_\sigma
&\longrightarrow \deleted{\mathcal G_{\rho,\sigma}}
 G\times G\times\unsure{\mathbb E^1|\mathbb E'}\times\unsure{\mathbb E^1|\mathbb E'}\\
(\underline u,\underline a,\underline b)
&\longmapsto
(\alpha=\underline u\underline a^{-1},\beta=\underline u\underline b^{-1},
\mu=\det\underline u,\varepsilon=\varepsilon_\sigma(\underline b)),
\end{aligned}\right.
\]

je dis que l'image est dans $\mathcal G^*_{\rho,\sigma}$, i.e.
qu'on aura bien
\[
\tag{\unsure{*}}
\alpha\rho(\alpha)^{-1}
=\varepsilon\bigl(\beta\sigma(\beta)^{-1}\bigr)\varepsilon_1(\mu-1).
\]
En effet
\[
\begin{aligned}
\alpha\rho(\alpha)^{-1}&=\underline u\rho(\underline u)^{-1}\\
\beta\rho(\beta)^{-1}
&=\underline u\underbrace{
\underline b^{\inserted{-1}}\sigma(\underline b^{-1})^{-1}
}_{\varepsilon_\sigma(\underline b)^{-1}=\varepsilon_\sigma(\underline b)}
\sigma(\underline u)^{-1}
=\varepsilon\underline u\sigma(\underline u)^{-1}
\end{aligned}
\]
et on sait (par II) que
\[
\underline u\rho(\underline u)^{-1}
=\underline u\sigma(\underline u)^{-1}\varepsilon_1(\mu-1)
\]
d'où $\unsure{(*)}$, puisque $\varepsilon^2=1$. Mais de
plus le morphisme (25) est mono, \deleted{donc}
car en termes de $(\alpha,\beta,\mu,\varepsilon)$, $u$ est déterminé
\inserted{(en termes de $\alpha$ déjà)},
(comme l'\uline{unique} section de $B'_0$ qui
soit de la forme $\alpha\underline a^{-1}$, avec $\underline a\in\Gamma T_\rho$
(car $T_\rho\cap B'_0=\{1\}$)), et $\underline b$ est alors déterminé
par la relation $\beta=\underline u\underline b^{-1}$. On trouve
donc que (25) est un \uline{iso}
\[
\tag{26}
\left\{\begin{aligned}
B'_0\times A_\rho\times N'_\sigma
&\overset{\sim}{\longrightarrow}\mathcal G^*_{\rho,\sigma}\\
(\underline u,\underline a,\underline b)
&\longmapsto
(\alpha=\underline u\underline a^{-1},\beta=\underline u\underline b^{-1},
\mu=\det\underline u,\varepsilon=\varepsilon_\sigma(\underline b))
\end{aligned}\right.
\]

\scanpage{318}
\mspage{459}

Si on a deux sections $(\underline u,\underline a,\underline b)$ et
$(\underline u',\underline a',\underline b')$ de
$B'_0\times T_\rho\times N'_\sigma$, à quelles
conditions a-t-on
\margin{\unsure{(27)}}
\[
\tag{27}
\widetilde\alpha=\widetilde\alpha',\quad
\widetilde\beta=\widetilde\beta'\ ?
\]
Il suffit évidemment qu'on ait
\[
\tag{28}
\underline a'=\lambda_1\underline a,\qquad
\underline b'=\lambda_2\underline b
\qquad \lambda_1,\lambda_2\in\Gamma G_{\mathrm m}
\]
et je dis que ces conditions sont aussi nécessaires.
Cela provient du fait que
$B'_0\hookrightarrow\unsure{PG}\simeq\unsure{GP(1)}$
est mono (soit $\widetilde B'_0$ l'image de $B'_0$),
que \deleted{\unsure{les}} $A_\rho\to PG$ et
$N'_\sigma\longrightarrow PG$ ont comme noyau $G_{\mathrm m}$
(\unsure{soit|soient} $\widetilde T_\rho$ et $\widetilde N'_\sigma$ leurs images),
et que l'on a dans $PG$
\[
\tag{29}
\widetilde T_\rho\cap\widetilde B'_0=\{1\},
\]
de sorte que \deleted{\illegible{1}} $\widetilde{\underline u}$
\inserted{$\in\Gamma\widetilde B'_0$ et $\widetilde{\underline a}\in\Gamma\widetilde A_\rho$}
sont définis \uline{de façon unique} en termes de $\widetilde\alpha$
par la relation
\[
\deleted{\illegible{1}}\qquad
\widetilde\alpha=\widetilde{\underline u}\widetilde{\underline a}^{-1}
\]
-- et qu'alors $\widetilde{\underline b}$ est défini
\inserted{en termes de $\widetilde\alpha,\widetilde\beta$}
par la relation
\[
\widetilde\beta=\widetilde{\underline u}\widetilde{\underline b}^{-1}.
\]

Ceci montre que \deleted{\unsure{\(\mathcal G\)}} l'on a un \uline{isomorphisme}
\[
\tag{30}
\left\{\begin{aligned}
\widetilde B'_0\times\widetilde A_\rho\times\widetilde N'_\sigma
&\overset{\sim}{\longrightarrow}
\widetilde{\mathcal G}^{*}_{\rho,\sigma}
\subset PG\times PG\times\unsure{\mathbb E^1|\mathbb E'}\\
(\widetilde{\underline u},\widetilde{\underline a},\widetilde{\underline b})
&\longmapsto
(\widetilde{\underline u}\widetilde{\underline a}^{-1},
\widetilde{\underline u}\widetilde{\underline b}^{-1},
\delta(\widetilde{\underline u}))
\end{aligned}\right.
\]

\scanpage{319}

où
\[
\tag{31}
\delta(\widetilde{\underline u})=\det(\underline u)
\]
quand $\widetilde{\underline u}$ provient de $\underline u\in\Gamma B'_0$ --
ou, ce qui revient au $\overline{\mathrm m}$ (cf p. 441)
$\delta(\widetilde{\underline u})$ est défini par
\[
\tag{32}
\widetilde{\underline u}(N_0)=\delta(\widetilde{\underline u})N_0.
\]

\margin{\uline{NB} Le \unsure{centralisateur} de $\widetilde\sigma$ dans
$\widetilde G=PG$ \uline{n'est pas} $\widetilde T_\sigma$, mais
$\widetilde N'_\sigma$ \deleted{\illegible{2}}.}

\uline{Rmq.} Notons qu'une section $(\widetilde\alpha,\widetilde\beta,\mu)$
\deleted{\unsure{$\in\Gamma$}} de $\widetilde{\mathcal G}^{*}_{\rho,\sigma}$
\unsure{reste} \deleted{\unsure{entièrement}} déterminée,
\inserted{\uline{mod op. de} $\unsure{\widetilde N'_\sigma}$ \uline{sur} $\beta$},
\unsure{pour} \unsure{être} \unsure{précis}, en termes de $\widetilde\alpha$ --
\unsure{condition|conditions} à satisfaire sur $\widetilde\alpha$, pour
qu'il existe $\widetilde\beta,\mu$ \deleted{\illegible{1}} qui la
\unsure{complètent} en une section de $\widetilde{\mathcal G}^{*}_{\rho,\sigma}$,
étant la suivante : le coefficient $D$ de
$\zeta=\alpha\rho(\alpha)^{-1}=\begin{pmatrix}A&B\\C&D\end{pmatrix}$
est inversible (ce qui permet de déterminer $\mu$ par $\mu=1/D$), puis
on détermine $\widetilde\eta$ \deleted{\illegible{1}} par
$\widetilde\eta=\widetilde\zeta\widetilde\varepsilon_1(\mu-1)$, et
 enfin on trouve \inserted{les} $\widetilde\beta$ \inserted{correspondants}
par la condition
\[
\widetilde\beta\widetilde\sigma(\widetilde\beta)^{-1}=\widetilde\eta.
\]

\margin{\uline{NB} $\widetilde\alpha$ dépend de trois paramètres,
tandis que $\widetilde{\mathcal G}^{*}_{\rho,\sigma}$
\uline{\unsure{est de|admet}} \unsure{dim.} relative $4=2+1+1$
(car $\widetilde{\mathcal G}^{*}_{\rho,\sigma}
\simeq\widetilde B'_0\times\widetilde T_\rho\times\widetilde N'_\sigma$).}

\margin{\uline{NB} On a relevé \unsure{loc$^{\mathrm t}$}
$\widetilde\alpha$ en $\alpha\in\Gamma G$,
$\zeta$ \unsure{ne} dépend du choix du relèvement.}

De $\overline{\mathrm m}$, une section $(\alpha,\beta,\mu,\varepsilon)$
\inserted{de $\mathcal G^*_{\rho,\sigma}$}
est déterminée par \deleted{\unsure{$\alpha$ seul}}
\inserted{\deleted{\illegible{2}} $\alpha$ seul} mod
\deleted{\illegible{2}} $N'_\sigma$ -- on trouve d'abord
$\underline u,\underline a$ en termes

\scanpage{320}
\mspage{460}

de $\alpha$ comme dessus, d'où $\mu=\det\underline u$, d'où
et il \deleted{suffit} \deleted{$\varepsilon\beta\sigma(\beta)^{-1}=\unsure{\alpha|\eta}$}
reste à déterminer $\beta$ comme
$\beta=\underline u\underline b^{-1}$, avec $\underline b\in\Gamma N'_\sigma$
arbitraire, $\varepsilon$ étant défini alors en termes du choix de $\underline b$
comme étant $\varepsilon_\sigma(\underline b)$. La condition pour que
$\unsure{\alpha_0|\alpha}$ fixé puisse se compléter par $(\beta,\mu,\varepsilon)$
en $(\alpha,\beta,\mu,\varepsilon)\in\Gamma\mathcal G^{\unsure{\alpha|*}}_{\rho,\sigma}$,
est tjrs la condition que le coeff. $D$ \deleted{de $\zeta=\alpha\unsure{+}$}
\[
\tag{33}
\zeta=\alpha\rho(\alpha)^{-1}=\begin{pmatrix}A&B\\C&D\end{pmatrix}
\]
soit inversible, ce qui \inserted{permet de} \unsure{déterminera|déterminer} $\mu$ par
\[
\tag{34}
\mu=1/D.
\]

Posons
\[
\tag{35}
\alpha=\begin{pmatrix}a&b\\c&d\end{pmatrix}
\]
et explicitons $\alpha\rho(\alpha)^{-1}$; avec
$\rho=\begin{pmatrix}0&1\\v&1\end{pmatrix}$

\margin{$\delta=\det\alpha=ad-bc$}
\[
\begin{aligned}
\alpha\rho(\alpha)^{-1}
&=\begin{pmatrix}a&b\\c&d\end{pmatrix}
\begin{pmatrix}0&1\\v&\deleted{\unsure{w}}\inserted{1}\end{pmatrix}
\frac1\delta\begin{pmatrix}d&-b\\-c&a\end{pmatrix}
\left(-\frac1v\right)
\begin{pmatrix}\deleted{\unsure{w}}\inserted{1}&-1\\-v&0\end{pmatrix}\\
&=-\frac1{\unsure{\delta}v}
\begin{pmatrix}bv&a+b\deleted{\unsure{w}}\\dv&c+d\deleted{\unsure{w}}\end{pmatrix}
\begin{pmatrix}d\deleted{\unsure{w}}+bv&-d\\-c\deleted{\unsure{w}}-av&c\end{pmatrix},
\end{aligned}
\]
d'où
\[
\deleted{\unsure{=\frac1{\delta v}}}
\]
\[
\tag{36}
\left\{\begin{aligned}
(-\delta v)A&=\unsure{v}b(d+\unsure{vb|bv})\unsure{-}(a+b)(c+va)
&(-\delta v)B&=-vbd+c(a+b)\\
(-\delta v)C&=vd(d+vb)-(c+d)(c+va)
&(-\delta v)D&=-vd^2+cd+c^2,
\end{aligned}\right.
\]
la quantité la plus essentielle étant donc
\[
\tag{37}
\left\{\begin{aligned}
(-\unsure{\delta}v)D&=-vd^2+cd+c^2=\det(d\rho+c),\\
\text{\deleted{\unsure{et après cela|et après calcul}} où }&
\left\{\begin{aligned}
\deleted{\unsure{\alpha}}\ \delta&=\det\alpha=ad-bc,\\
v&\ \text{le \unsure{paramètre|coefficient} qui définit }
\rho=\begin{pmatrix}0&1\\v&1\end{pmatrix}
\end{aligned}\right.
\end{aligned}\right.
\]
